\documentclass[11pt,a4paper]{ctexart}
\usepackage[margin=2.4cm]{geometry}
\usepackage{amsmath,amssymb}
\usepackage{fancyhdr}
\usepackage{enumitem}
\linespread{1.18}
\setlength{\headheight}{15pt}
\setlength{\parskip}{0pt}

\pagestyle{fancy}
\fancyhf{}
\fancyhead[L]{\small Team 2216547}
\fancyhead[R]{\small \thepage}
\renewcommand{\headrulewidth}{0.4pt}
\fancypagestyle{plain}{\fancyhf{}\renewcommand{\headrulewidth}{0pt}}

\ctexset{
  section/format = \Large\bfseries\raggedright,
  subsection/format = \large\bfseries\raggedright,
  subsubsection/format = \normalsize\bfseries\raggedright
}
\setcounter{secnumdepth}{3}

\begin{document}

%====================== 摘要页 ======================
\thispagestyle{empty}
\begin{center}
{\small 2022 MCM/ICM\quad 选题 A\quad 队伍控制号 2216547\quad 摘要页（Summary Sheet）}\\[18pt]
{\LARGE\bfseries 基于生理学、动力学与气象学的最优骑行方案}\\[8pt]
{\large （Optimal Cycling Solution Based on Physiology, Dynamics and Meteorology）}\\[10pt]
\end{center}

\begin{center}
{\large\bfseries 摘\quad 要}
\end{center}

\vspace{6pt}
本文旨在刻画不同骑手的功率曲线，并研究骑手在赛道进程中应当如何分配自身的功率输出，才能在个人计时赛（Individual Time Trial）中取得更佳的运动表现。

首先，我们设计了肌肉疲劳—恢复模型（Muscle Fatigue-Recovery Model）与能量储备模型（Energy Reserve Model），并将二者耦合，得到骑手功率输出模型（Rider Power Output Model）。随后，我们利用文献[4][5]中“最大努力”（maximum-effort）实验的数据对模型参数进行校正。在此基础上，我们以男性与女性的计时赛专家（Time Trial Specialist，TTS）以及冲刺手（Sprinter）为例，通过所建立的模型得到了各自的功率曲线。

其次，我们从谷歌地球（Google Earth）采集真实道路点的坐标信息，构建道路的数据结构并定义其属性字段，从而在计算机中重新表达道路的高程、长度、曲率等属性。由此，自行车赛道得以可视化。接着，利用赛道信息，我们建立了自行车的牛顿力学模型，并基于牛顿第二定律建立了骑手的运动方程。

随后，在上述三个模型的基础上，我们形式化地定义了功率规划的优化问题。该问题的最优解由动态规划（Dynamic Programming）算法求得。进一步地，在该模型之上，我们引入了风速、风向与地面湿度等环境因素。此外，我们还考察了骑手自身对方案执行误差所造成的影响。

最后，我们修改前述模型的参数以模拟团队计时赛（Team Time Trial）。在团队计时赛中，车队常采用“山地火车”（mountain train）战术：队员排成一列，以尽可能降低整体空气阻力；相关信息的获取与数据处理参见文献[13]。

基于上述建模，我们可以得出以下结论：

\begin{enumerate}
\item 对于一段起伏相对明显的道路，冲刺手的最优功率规划方案是在起步阶段与上坡路段提高功率输出，而在相对平缓的路段以较低的功率骑行以保存体能。对于计时赛专家而言，情况恰好相反。
\item 性别对上述策略没有显著影响。
\item 地面湿度越高、逆风越强，骑手完成比赛所需的时间越长；而地面湿度越低、顺风越强，则用时越短。
\item 若骑手未严格遵循高度细化的方案、未能命中目标功率，其完赛时间将会延长。
\item 在团队计时赛中，得益于团队比赛所采用的“山地火车”策略，每名骑手所承受的平均空气阻力下降，使车队能够以更快的节奏推进。
\end{enumerate}

\vspace{6pt}
\noindent\textbf{关键词：}功率曲线；牛顿力学；动态规划；自行车运动

\newpage
\pagestyle{fancy}

%%====================== 正文 ======================
\section{引言}

\subsection{问题背景}

自行车赛事在世界范围内广受欢迎，其中尤以个人计时赛（Individual Time Trials，ITTs）与团队计时赛（Team Time Trials，TTTs）为甚。个人计时赛又被称为“真理之赛”（the race of truth），因为胜负既取决于每名骑手的力量与耐力，也取决于骑手依据道路状况分配功率的水平。为使骑手取得更好的成绩，通过数学建模来分析具有不同身体条件（体能储备、身体恢复能力、爆发能力等）的骑手应如何在不同的道路条件下科学、合理地骑行，具有十分重要的意义。

\subsection{问题重述}

\begin{enumerate}
\item 首先，要求我们依据骑手的身体属性建立关于骑手的模型，该模型能够输出骑手的功率曲线，用以表明骑手在给定功率水平下能够维持多长时间。
\item 在对骑手本身完成建模之后，我们需要对骑手在通过给定计时赛赛道时应如何施加功率进行建模。为解决该问题，我们应当考虑道路条件、环境因素以及骑手自身的模型。
\item 在此之后，我们应当考虑当骑手的实际表现偏离我们的模型时所得到的结果变化。
\item 最后，我们需要在模型中求解团队计时赛情形。
\end{enumerate}

为求解上述问题，我们按如下步骤推进：

\begin{itemize}
\item 设计骑手功率输出模型。我们设计骑手功率输出模型，并对该模型的参数进行校正。
\item 对道路进行离散化与可视化。我们从谷歌地球获取道路数据，将其离散为路段，并借助 Python 将其可视化。
\item 设计骑行的牛顿力学。通过分析自行车的物理模型，并将其与赛道的两点环境数据相结合，我们提出骑行的牛顿力学。
\item 使用动态规划求解问题。我们将自行车计时赛的提速策略视为一个最优控制问题，并使用动态规划求解该最优控制问题。
\item 修改模型参数以适配不同问题。我们修改模型参数以模拟运动员在不同环境条件下的表现；同样地，我们修改模型参数以模拟团队比赛中的团队表现。
\end{itemize}

\subsection{文献综述}

数十年来，科学家为理解骑手运动表现背后的生理学概念开展了大量工作。其中，临界功率（critical power，CP）与无氧工作能力（anaerobic work capacity，AWC）这两个概念在该主题近期的文章中屡被讨论，也有文章基于 CP 与 AWC 提出了更为经典的肌肉疲劳与恢复模型[1]。不仅如此，一些文章提出了基于静力推导的肌肉疲劳恢复模型[2]。此外，基于自行车运动员身体表现的功率输出模型也具有一定的有效性[3]。

科学家们不仅在人体肌肉模型方面开展工作，在人体数据测量领域亦有建树。Dajo Sanders 等人使用了世界自行车巡回赛车队中职业自行车手的数据，他们以最大有氧功率、峰值冲刺功率输出以及描述功率随时间衰减的指数常数，为每位受试者建立了功率—持续时间关系[4]。N. J. Rehrer 等人测量了运动员在为期六天、共十个赛段的自行车比赛中的 EI（能量摄入）与 EE（能量消耗）[3]。Tammie R. Ebert 等人采集了 31 名澳大利亚国家男子公路自行车手在 6 年、207 场比赛中的功率输出数据，并在实验室中对受试者进行最大递增负荷运动测试，以确定其最大有氧功率输出[5]。

S. Alireza Fayazi 等人将自行车计时赛的提速策略描述为一个最优控制问题，并使用动态规划求解该最优控制问题[2]。动态规划既是一种数学优化方法，也是一种计算机编程方法。该方法由 Richard Bellman 于 20 世纪 50 年代提出。动态规划指的是以递归的方式将复杂问题分解为更简单的子问题，从而简化复杂问题[6]。

\section{假设与合理性说明}

为简化问题并便于我们模拟真实条件，我们作出如下假设，并对每一条假设给出合理说明。

\begin{itemize}
\item \textbf{我们所采集的全部数据都是准确的。}我们通过各种途径采集的数据是在不同环境、不同时间、不同赛道上测量得到的，难免受到未被纳入考虑的变量所引起的误差干扰；同时，用于测量这些数据的设备也不能完全避免测量误差。因此在我们的模型中，我们假设所采集的全部数据都是准确的。
\item \textbf{在比赛开始时，运动员总是达到其最佳的初始身体状态。}由于比赛前数日骑手的睡眠与饮食会带来影响，骑手在比赛开始时的初始身体状态未必总是其所能达到的最佳初始状态。这种不确定性对建模而言颇为棘手。在我们的模型中，我们假设在比赛开始时，运动员总是达到其最佳的初始身体状态。
\item \textbf{运动员总是为当前道路选择最合适的自行车挡位。}在平路上选择较大的挡位效率低下，而在坡道上选择较小的挡位既费时又费力。更为合理的做法是在平路上选择较小的挡位、在坡道上选择较大的挡位。在我们的模型中，我们假设骑手在不同路段上总是能够选择相应最可用的挡位。
\item \textbf{运动员比赛过程中不出现意外情况。}我们假设所模拟的运动员在比赛中不会身体不适、跌倒或受伤。我们假设运动员是在理想条件下完成比赛的。
\item \textbf{环境参数在比赛过程中保持恒定。}部分自行车比赛会持续数小时乃至数天，其间天气与环境可能发生显著变化，例如由雨转晴。但在我们的模型中，我们假设比赛期间环境参数恒定。
\item \textbf{男、女骑手完成相同的圈数。}为了比较男、女骑手在同一赛道上的表现，我们假设男、女骑手完成相同的圈数。
\end{itemize}

\section{符号说明}

本节定义模型中的主要符号，各符号的具体取值将在后文给出。

\begin{itemize}
\item $P_{max}$：当前状态下肌肉所能释放的最大功率，单位 W；
\item $P_{max0}$：初始状态下肌肉所能释放的最大功率，单位 W；
\item $P_{rider}$：骑手在当前时刻的输出功率，单位 W；
\item $k_f$：肌肉疲劳系数，无量纲；
\item $k_r$：肌肉恢复系数，无量纲；
\item $E$：当前时刻的身体能量储备，单位 J；
\item $E_0$：初始身体能量储备，单位 J；
\item $P_{in}$：能量恢复速率，单位 W；
\item $P_{out}$：能量输出速率，单位 W；
\item $k_{in}$：能量恢复系数，无量纲；
\item $m_b,\ m_r,\ m_t$：自行车质量、骑手质量以及二者总质量，单位 kg；
\item $k_{wheel}$：自行车功率转换系数，无量纲；
\item $A$：骑手的迎风面积，单位 $\mathrm{m^2}$；
\item $C_d$：空气阻力系数，无量纲；
\item $v$：骑手的速度，单位 $\mathrm{m/s}$；
\item $\rho$：空气密度，单位 $\mathrm{kg/m^3}$；
\item $\mu$：滚动摩擦系数，无量纲；
\item $\theta$：道路的倾角，单位 rad。
\end{itemize}

\section{我们的模型}

\subsection{骑手功率输出模型}

自行车手在比赛中的功率输出由多种因素决定，例如运动员的睡眠、饮食、情绪、肌肉状态（最大功率输出、抗疲劳能力、疲劳恢复能力等）、体能状态（初始体能、体能恢复能力）、比赛时的天气与气温等。因此，对于这样一个复杂的人体功率输出系统，我们以两个单一模型相耦合的方式来考察：肌肉疲劳—恢复模型与人体能量储备模型。我们用真实运动员的数据校正模型的参数，最终得到不同运动员的功率曲线。由我们的模型得到的运动员功率曲线与真实运动员的功率曲线吻合良好。

\subsubsection{肌肉疲劳—恢复模型}

自行车手主要的产力肌为股四头肌、臀大肌与臀中肌。踩踏自行车所能施加的最大功率是肌肉疲劳程度的递减函数。本文采用单状态动态模型来刻画运动员当前肌肉所能输出的最大功率的变化。文献[7]的结果为“将促成疲劳的主要生理因素之影响加以整合”提供了依据。

我们假设 $P_{max}$ 为运动员肌肉在当前状态下所能输出的最大功率，其初始值为 $P_{max0}$（即运动员在肌肉最佳状态下所能输出的最大功率数值），$P_{rider}$ 为运动员肌肉当前的功率输出。$P_{rider}$ 越大，运动员的肌肉疲劳越快，$P_{max}$ 下降也越快。文献[8]中若干实验的结果支持如下模型：

\begin{equation}
\left.\frac{\mathrm{d}P_{max}(t)}{\mathrm{d}t}\right|_{\text{fatigue}}
= -k_f\cdot\frac{P_{max}(t)\cdot P_{rider}(t)}{P_{max0}}
\end{equation}

其中 $k_f$ 为肌肉抗疲劳系数，$k_f$ 因人而异。肌肉恢复也几乎呈指数增长，肌肉恢复过程可由如下模型表示[8]：

\begin{equation}
\left.\frac{\mathrm{d}P_{max}(t)}{\mathrm{d}t}\right|_{\text{recovery}}
= k_r\cdot\left(P_{max0}-P_{max}(t)\right)
\end{equation}

其中 $k_r$ 为肌肉疲劳恢复系数，$k_r$ 因人而异。

然而，运动员肌肉的疲劳与恢复在时间上并非完全分离。每当一组肌纤维处于激活状态时，其中一部分会疲劳，另一部分则会恢复[2]。换言之，肌肉的疲劳与恢复是同时发生的。为反映这种同时性，我们将式（1）与式（2）合并，得到肌肉疲劳—恢复模型：

\begin{equation}
\frac{\mathrm{d}P_{max}(t)}{\mathrm{d}t}
= -k_f\cdot\frac{P_{max}(t)\cdot P_{rider}(t)}{P_{max0}}
+ k_r\cdot\left(P_{max0}-P_{max}(t)\right)
\end{equation}

\subsubsection{能量储备模型}

我们知道，在人体运动过程中，血糖水平会在一定范围内波动。血糖是有氧运动（如自行车骑行）的主要能量来源。随着血糖被消耗，肝糖原被分解为血糖以提高血糖水平。同时，自行车比赛不会耗尽运动员的肝糖原，这是常识。因此，血糖浓度在比赛过程中能够稳定在一定范围内。据此，我们设计了如下能量储备模型：

\begin{equation}
P_{in}(t)=k_{in}\cdot\left(P_{max0}-P_{max}(t)\right)
\end{equation}

\begin{equation}
P_{out}(t)=P_{rider}(t)+P_{other}(t)=k_{out}\cdot P_{rider}(t)
\end{equation}

\begin{equation}
E(t)=E_0+\int_{0}^{t}P_{in}(t)\,\mathrm{d}t-\int_{0}^{t}P_{out}(t)\,\mathrm{d}t
\end{equation}

其中 $E$ 为可直接利用的能量总量，可用以表示血糖浓度；$E_0$ 为其初始值，即运动员经充分休息后可直接利用的能量总量；$P_{in}$ 为能量补充速率，可用以表示肝糖原分解的速率；$P_{out}$ 为能量输出速率，可用以表示血糖消耗的速率；$P_{other}$ 为运动员除机械输出之外的生理活动所消耗能量的功率。

式（4）的含义是：能量补充速率与肌肉疲劳程度成正比，其中 $k_{in}$ 为能量补充系数。

式（5）的含义是：运动员的能量输出速率由肌肉的机械输出速率与其他生理活动的能量消耗速率组成。在整个比赛过程中，$P_{out}$ 近似正比于 $P_{rider}$，$k_{out}$ 即为该比例参数。

运动员的输出功率同时由肌肉疲劳—恢复模型与能量储备模型决定。

\begin{enumerate}
\item 当运动员肌肉疲劳而能量储备充足时，运动员肌肉的最大输出功率受肌肉疲劳程度的限制。
\item 当运动员肌肉的能量储备耗尽而肌肉尚未过度疲劳时，运动员肌肉的最大输出功率受能量储备的限制。
\end{enumerate}

肌肉疲劳—恢复模型与能量储备模型共同作用的效果如下。能量储备较为充裕的运动员，其功率曲线表现为：最大功率输出仅受肌肉疲劳程度的限制，曲线上不出现转折点；能量储备处于平均水平的运动员，其功率曲线在时间较短的一段区间内由肌肉疲劳程度决定最大功率输出；能量储备偏低的运动员，其肌肉输出速率受能量储备限制的时间区间比前者更短，即曲线出现转折的位置更靠前。

\subsubsection{参数校正}

在我们的模型中，一名运动员的能力可以完整地由六个参数表示：$k_f$、$P_{max0}$、$k_r$、$E_0$、$k_{in}$、$k_{out}$。为针对不同运动员获得不同的模型表示，我们为不同运动员设计不同的参数。我们使用文献[4]中“最大努力”（Maximum-Effort）实验的数据来设计参数：“最大努力”实验指的是运动员经充分休息后在平坦路段以全力冲刺，并记录该过程中运动员的功率输出。我们求出了与真实运动员“最大努力”实验数据拟合最佳的参数，记录如下。

对于不同类型的运动员，参数取值如下：

\begin{itemize}
\item 计时赛专家（男）：$k_f=0.0416$，$k_r=0.0104$，$k_{in}=1.11$，$k_{out}=1.50$，$P_{max0}=1000$，$E_0=10^{6}$；
\item 计时赛专家（女）：$k_f=0.018$，$k_r=0.00901$，$k_{in}=1.18$，$k_{out}=1.46$，$P_{max0}=900$，$E_0=8\times10^{5}$；
\item 冲刺手（男）：$k_f=0.057$，$k_r=0.001078$，$k_{in}=0.50$，$k_{out}=1.57$，$P_{max0}=1500$，$E_0=10^{6}$；
\item 冲刺手（女）：$k_f=0.046$，$k_r=0.0009$，$k_{in}=0.53$，$k_{out}=1.52$，$P_{max0}=1300$，$E_0=8\times10^{5}$。
\end{itemize}

\subsubsection{功率曲线}

功率曲线表示骑手在给定持续时间长度内所能维持的最大功率。我们采用如下算法得到不同类型运动员的功率曲线。

算法思想：固定输出功率 $P$，求运动员能够坚持该功率输出的最长时间长度。

\noindent\textbf{算法 1}\quad 绘制功率曲线

\noindent\textbf{输入：}$k_f,\ k_r,\ k_{in},\ k_{out},\ P_{max0},\ E_0$，即该运动员的参数。

\noindent\textbf{输出：}该运动员的功率曲线。

\begin{enumerate}
\item 初始化列表 $t_{list}=[\ ]$，$P_{rider\_list}=[\ ]$；
\item 令 $P_{rider}$ 从 100 递增至 1500，步长为 1，对每个 $P_{rider}$ 执行第 3 至第 6 步；
\item 令 $E=E_0$；
\item 令 $i$ 从 1 递增至 799，在每一步中更新 $P_{max}$ 的取值，并更新 $E$ 的取值；若 $P_{rider}>P_{max}$ 或 $E<0$，则令 $t=i$ 并跳出该层循环；
\item 将 $t$ 追加到 $t_{list}$，将 $P_{rider}$ 追加到 $P_{rider\_list}$；
\item 以 $t_{list}$ 为横坐标、$P_{rider\_list}$ 为纵坐标绘图，即得该运动员的功率曲线。
\end{enumerate}

\subsection{道路离散化与虚拟化}

为将我们的模型应用于各类计时赛赛道，我们利用真实世界中两条自行车赛道的数据构建了离散化的虚拟赛道：日本东京的 2021 年奥运会计时赛赛道，以及比利时佛兰德的 2021 年 UCI 世界锦标赛计时赛赛道；此外，我们还自行设计了一条合理的自行车赛道。

\subsubsection{赛道的离散化}

我们在谷歌地球上标定真实自行车赛道上的离散数据点，每个点包含其纬度、经度与海拔信息。对于直线道路，构成该段道路的两个离散数据点间距较大；而对于弯道，离散数据点的间距更小，以便获得更高的精度。

对于从谷歌地球获取的一个点，我们将其定义为一个三维向量
\[
p'=(x',y',z')
\]
其中 $x'$ 为纬度，$y'$ 为经度，$z'$ 为高度。

由于信息点的前两个维度存储的是纬度与经度信息，我们将其转换为以米为单位的平面二维坐标，以获得更高的精度。于是，我们从赛道上选取一个参考点作为原点，并依据地球半径计算出前两个维度相对于原点的长度，从而使我们的每个点变为
\[
p=(x,y,z)
\]
其中 $x$ 与 $y$ 为该点在其赛道地图上相对于原点的坐标长度，$z$ 保持不变。这样，我们便可以在程序仿真中在坐标轴上重建并显示整幅地图。

\subsubsection{道路的数据结构定义}

在从谷歌地球采集到点的信息之后，我们定义道路的数据结构。在我们的模型中，道路基本上由两个点构成。

对于道路 $road_i$，其中 $i$ 为该道路在道路列表中的索引，假设点 $p_i^{(1)}$ 与 $p_i^{(2)}$ 分别为 $road_i$ 的起点与终点，则我们可以得到如下属性：

\begin{equation}
\begin{cases}
L_i=\sqrt{\left(p_i^{(1)}.x-p_i^{(2)}.x\right)^2+\left(p_i^{(1)}.y-p_i^{(2)}.y\right)^2+\left(p_i^{(1)}.z-p_i^{(2)}.z\right)^2}\\[0.6em]
\Delta h_i=p_i^{(1)}.z-p_i^{(2)}.z\\[0.4em]
\theta_i=\arctan\dfrac{\Delta h_i}{\sqrt{\left(p_i^{(1)}.x-p_i^{(2)}.x\right)^2+\left(p_i^{(1)}.y-p_i^{(2)}.y\right)^2}}
\end{cases}
\end{equation}

其中 $L_i$ 为 $road_i$ 的长度，$\Delta h_i$ 为 $road_i$ 的高差，$\theta_i$ 为 $road_i$ 的倾角。

\subsubsection{点与道路的曲率半径定义}

为使模型能够纳入弯道的影响，我们引入点与道路的曲率半径概念。

对于赛道上的一个给定数据点 $p_i$，其曲率半径记为 $R(p_i)$，即由 $p_{i-1}$、$p_i$ 与 $p_{i+1}$ 三点所构成三角形的外接圆半径。若 $p_{i+1}$ 或 $p_{i-1}$ 不存在，即 $p_i$ 为整条赛道的端点之一，则 $p_i$ 的曲率半径不予定义。

对于一条给定道路，其端点为 $p_i$ 与 $p_{i+1}$；若两端点的曲率半径均存在，则该道路的曲率半径为 $\dfrac{R(p_i)+R(p_{i+1})}{2}$。否则，该道路为赛道的起始路段或终止路段：起始路段取其终点 $p_{i+1}$ 的曲率半径，终止路段取其起点 $p_i$ 的曲率半径。

曲率半径越大，该段道路越接近直线。

\subsection{骑行的牛顿力学分析}

\subsubsection{直线赛道上的骑行}

根据牛顿第二定律并忽略车轮转动的惯性效应，自行车的纵向模型假设为

\begin{equation}
m_t a=k_{rider}\frac{P_{rider}}{v}-\frac{1}{2}C_d\rho A v^2-m_t g\left(\mu\cos\theta+\sin\theta\right)-F_b
\end{equation}

其中 $m_t$ 为骑手与自行车的总质量，$a$ 表示骑手的加速度。因此，依据牛顿第二定律，$m_t a$ 等于施加于该系统的合力。

一般而言，自行车手在规划自身行进时有两个自由度：规划自身输出功率的自由度，以及依据道路信息规划自身挡位组合的自由度。我们假设自行车手均为训练有素的选手，能够基于经验最优地选择挡位组合。$k_{rider}$ 表示在选手所选挡位组合下，选手输出功率向自行车的转换效率。因此，$k_{rider}\dfrac{P_{rider}}{v}$ 表示自行车手向自行车输出的力。

$\dfrac{1}{2}C_d\rho A v^2$ 表示自行车手所受空气阻力的大小，其中 $C_d$ 为空气阻力系数，$\rho$ 为赛道上空气的密度（其随海拔升高而减小），$A$ 与 $v$ 分别为骑手的迎风面积与速度。

$m_t g\left(\mu\cos\theta+\sin\theta\right)$ 一项是摩擦力 $\mu m_t g\cos\theta$ 与骑手系统所受下滑分力 $m_t g\sin\theta$ 的合力。$\mu$ 为自行车轮胎与道路之间的滚动摩擦系数，$\theta$ 为给定道路的倾角。

$F_b$ 表示当选手决定刹车时刹车片对轮胎的阻力。

上述各变量的典型取值如下：自行车与骑手的总质量 $m_t=60\,\mathrm{kg}$；自行车功率转换系数 $k_{wheel}=0.99$；骑手迎风面积 $A=0.35\,\mathrm{m^2}$；空气阻力系数 $C_d=0.05$；空气密度 $\rho=1.1\,\mathrm{kg/m^3}$；滚动摩擦系数 $\mu=0.003$。

\subsubsection{弯道上的骑行}

当自行车手在弯道上骑行时，需要有一部分静摩擦力来提供转向所需的向心力。

对于一段曲率半径为 $R$ 的道路，当骑手速度为 $v$、质量为 $m_t$ 时，所需的向心力为 $\dfrac{m_t v^2}{R}$。因此，若最大静摩擦力为 $\mu_0 m_t g\cos\theta$，则骑手的速度不能超过 $\sqrt{\mu_0 g\cos\theta R}$。

\subsection{优化问题的形式化与数值求解}

在以上各节中，我们系统地建立了关于自行车手身体肌肉、自行车物理属性以及赛道物理属性的数学模型。在本节中，我们尝试形式化地定义个人计时赛中的最短时间优化问题，并使用动态规划（DP）求解该问题。

\subsubsection{问题的形式化}

在个人计时赛中，目标是以尽可能短的时间抵达终点。因此，我们将优化问题形式化定义如下。

\begin{equation}
\min_{P_{rider}(t)}T=\min_{P_{rider}(t)}\left\{\int_{t_0}^{t_f}\mathrm{d}t\right\}=\min_{P_{rider}(t)}\left\{\int_{x_0}^{x_f}\frac{\mathrm{d}x}{v(t)}\right\}
\end{equation}

约束条件为：

\begin{equation}
\begin{cases}
\text{速度限制：}\quad 0\leqslant v(t)\leqslant \min\left\{v_{max},\ \sqrt{\mu g\cos\theta R}\right\}\\[0.4em]
\text{剩余能量限制：}\quad 0\leqslant E(t)\leqslant E_0\\[0.4em]
\text{骑手功率限制：}\quad 0\leqslant P_{rider}(t)\leqslant P_{max}(t)
\end{cases}
\end{equation}

其中决策变量 $P_{rider}(t)$ 表示骑手在 $t$ 时刻决策的输出功率，$t_0$、$t_f$ 分别表示赛道起点与终点所对应的时刻，$x_0$、$x_f$ 分别表示赛道起点与终点的位置。

\subsubsection{用动态规划数值求解优化问题}

动态规划既是一种数学优化方法，也是一种计算机编程方法。动态规划的主要思想十分直接：我们将一个问题划分为更小的嵌套子问题，然后组合这些子问题的解以得到整体解。这一思想被称为最优性原理，本节给出更为正式的阐述。“动态规划”一词最早由 Richard Bellman 于 20 世纪 40 年代使用，用以描述需要逐一寻找最佳决策的问题；在 20 世纪 50 年代，他对其加以提炼，用以描述将小的决策问题嵌套进更大问题之中的情形。最优性原理的数学表述以他的名字被记为贝尔曼方程[6]。

为便于讨论，我们定义如下记号。

我们称四元组 $s_i=(x_i,v_i,P_{max,i},E_i)$ 为一个状态。令四个状态变量 $x_i$、$v_i$、$P_{max,i}$、$E_i$ 分别表示当前位置与起点的距离、当前状态的速度、当前状态的最大输出功率以及当前状态的剩余能量。

我们定义：状态 $s_j=(x_j,v_j,P_{max,j},E_j)$ 可以在代价 $\Delta t$ 下转移到状态 $s_i=(x_i,v_i,P_{max,i},E_i)$（其中 $x_j\leqslant x_i$），当且仅当存在 $P_{rider}$，$0\leqslant P_{rider}\leqslant P_{max,j}$，满足

\begin{equation}
\begin{cases}
\Delta t=\dfrac{x_i-x_j}{v_j}\\[0.6em]
v_i=v_j+a\Delta t\\[0.4em]
P_{max,i}=-k_f P_{max,j}\dfrac{P_{rider}}{P_{max0}}\Delta t+k_r\left(P_{max0}-P_{max,j}\right)\Delta t\\[0.6em]
E_i=E_j+k_{in}\left(P_{max0}-P_{max,j}\right)\Delta t-P_{rider}\Delta t
\end{cases}
\end{equation}

我们称 $P_{rider}$ 为从 $s_j$ 到 $s_i$ 的决策。

多次转移的代价为各次转移代价之和。令 $T(x_i,v_i,P_{max,i},E_i)$ 表示从起始状态 $(0,0,P_{max0},E_0)$ 最优地转移到当前状态 $(x_i,v_i,P_{max,i},E_i)$ 所需的最小代价。

依据贝尔曼最优性原理，即最优策略（决策集合）具有这样的性质：无论初始状态与初始决策如何，其余的决策必须相对于第一次决策所产生的状态构成一个最优策略。

贝尔曼最优性原理对我们的优化问题成立。因此我们有

\begin{equation}
T(x_i,v_i,P_{max,i},E_i)=\max_{P_{rider}}\left\{T(x_j,v_j,P_{max,j},E_j)+\Delta t\right\}
\end{equation}

其中 $P_{rider}$、$\Delta t$ 分别为从状态 $(x_j,v_j,P_{max,j},E_j)$ 转移到 $(x_i,v_i,P_{max,i},E_i)$ 的决策与代价。

\subsubsection{动态规划解法的复杂度分析}

由于现代电子计算机直接处理连续数据的能力十分有限，我们将连续数据量化为离散层级。

令 $n_x$、$n_v$、$n_p$、$n_e$、$n_r$ 分别表示 $x$、$v$、$P_{max}$、$E$、$P_{rider}$ 的量化级数（互不相同的离散层级的数目）。

因此，状态总数为 $n_x n_v n_p n_e$，而对每个状态所需的计算次数可表示为 $O(n_r)$。于是，我们动态规划算法的时间复杂度为 $O(n_x n_v n_p n_e n_r)$，空间复杂度为 $O(n_x n_v n_p n_e)$。

当 $n_x=400$、$n_v=30$、$n_p=10$、$n_e=10$、$n_r=450$ 时，在 MacBook Pro（2020 年款，Apple M1）上完成计算约需 10 分钟。
\section{结果}

\subsection{不同性别计时赛专家与冲刺手的功率档案}

依据前述算法与校正后的参数，我们得到计时赛专家与冲刺手在男、女两种情形下的功率档案，即各骑手在给定持续时间长度内所能维持的最大功率随持续时间变化的曲线。

\subsection{不同赛道的功率分配}

\subsubsection{日本东京的 2021 年奥运会计时赛赛道}

针对不同类型的骑手，动态规划算法给出了赛道各分段上的最优功率分配。

由功率分配曲线可以看出，冲刺手适合在比赛开始阶段以及上坡阶段提高功率输出，从而在比赛进程中取得领先。然而，与计时赛专家相比，冲刺手通常不擅长长距离骑行，因此在较长的平路路段上，他们往往落后于计时赛专家。

据此，我们提出如下规划策略：对于一段起伏相对明显的道路，冲刺手的最优功率规划方案大致是在起步阶段与大上坡路段提高功率输出，而在较长、相对平缓的路段以较低功率骑行以保存体能。对于计时赛专家而言，情况恰好相反：不在比赛开始时全力加速，上坡路段亦然，而是把更多的功率分配到较长且平缓的路段。性别对这一策略没有显著影响。

\subsubsection{2021 年 UCI 世界锦标赛计时赛赛道与本文设计的赛道}

我们分别给出了 2021 年 UCI 世界锦标赛计时赛赛道与我们所设计赛道上的最优功率分配。这两条赛道的数据同样支持我们关于计时赛专家与冲刺手的上述观点。

\subsection{结果分析}

我们给出了计时赛专家（男）在 2021 年奥运会计时赛赛道上的功率输出、速度与赛道高度分布。可以看到，当赛道为上坡时，该骑手的功率较大而速度较小；下坡时则相反。这进一步印证了我们模型的有效性。

\subsection{气象条件的潜在影响}

即便在同一条赛道上，不同的天气条件也会造成不同的比赛结果。在我们的模型中，我们主要考虑风与地面湿度对赛道表现的影响。

\subsubsection{风}

风对自行车比赛的影响因素可分为风向与风力。在顺风较强的情形下，运动员需要克服的空气阻力较小；在逆风较强的情形下，运动员需要克服的空气阻力较大。针对不同的风况，我们对模型作如下不同修正：

\begin{equation}
v=\left|v_{rider}\pm v_{wind}\right|\cdot\cos\theta
\end{equation}

其中 $v_{rider}$ 为自行车手的速度，$v_{wind}$ 为风速，$\theta$ 为骑手行进方向与风向之间的夹角。

\subsubsection{地面湿度}

若地面湿度较高，自行车手会降低速度以避免打滑或摔倒；而若地面干燥，自行车手则可以尽情释放速度。在我们的模型中：地面湿度越高，$\mu$ 越小；地面湿度越低，$\mu$ 越大。

环境因素对骑手表现的影响可由仿真结果得到，考察对象为计时赛专家（男）在同一赛道上的“距离—功率”曲线：

\begin{itemize}
\item 理想条件下的“距离—功率”曲线，完成时间为 24 分 25 秒；
\item 雨天条件下同一计时赛专家（男）的“距离—功率”曲线，完成时间为 30 分 10 秒；
\item 逆风条件下同一计时赛专家（男）比赛时的“距离—功率”曲线，完成时间为 25 分 11 秒。
\end{itemize}

由仿真结果可以看到，地面湿度越高、逆风越强，骑手完成比赛所需的时间越长。

\subsection{结果对骑手偏离目标功率分配的敏感性}

由于环境因素不可控、骑手自身身体条件存在差异，以及骑手与教练对赛道的看法与分段方式不同，骑手不会严格遵循所规划的最优方案进行比赛。我们将骑手的实际功率输出视为规划的最优功率输出方案与一个服从正态分布的随机误差之和。

\begin{equation}
P_{rider}(t)=P_{rider\_real}(t)=P_{rider\_plan}(t)+\varepsilon(t)
\end{equation}

其中 $\varepsilon(t)$ 为服从正态分布的随机误差。

将改进后的模型带入仿真，得到如下结果：最优情形下的完赛时间为 24 分 25 秒，受到扰动后的完赛时间为 26 分 44 秒。

\subsection{团队计时赛的结果}

我们知道，自行车手在高速骑行时需要克服空气阻力。速度越快，空气阻力越大。在团队骑行中，车队常采用“山地火车”战术：队员排成一列，以尽可能降低整体空气阻力。位于队列最前方的骑手需要承受最大的空气阻力，而位于队列中越靠后的骑手所需承受的空气阻力越小。通过资料查找与数据处理[13]，我们得到了 6 人车队中每名骑手所受空气阻力的信息，其中百分比表示每名骑手所需承受的空气阻力与队列最前方骑手相比的比例。

车队在“山地火车”战术中所采用的动态策略是：位于队列最前方的骑手每次为车队承担最大空气阻力一段时间，随后换到队列最后一位，享受最小的空气阻力；其余骑手依次前移、轮番顶替，如此循环往复。对于不同的替换间隔时长（即在每个位置轮转周期中一名骑手处于队列最前方的时间），车队中每名骑手的平均功率输出会不同，亦即车队的平均速度会不同。经过多次仿真，我们找到了使平均功率输出最大（完赛时间最短）的替换间隔时长。

仿真结果显示：团队计时赛的“距离—功率”曲线对应完赛时间为 22 分 36 秒，个人计时赛的“距离—功率”曲线对应完赛时间为 24 分 25 秒。由仿真结果可以看到，得益于团队比赛所采用的“山地火车”策略，每名骑手所承受的平均空气阻力下降，使车队能够以更快的节奏推进。

\section{优缺点}

\subsection{优点}

\begin{itemize}
\item 我们依据人体肌肉生理学与人体能量模型建立了骑手的功率输出模型，并通过考虑骑手不同的输入参数，输出带有各自属性字段的功率档案。
\item 我们以真实数据与合理设计对自行车赛道建模，所设定的数据结构包含高差、曲率半径等属性，并包含地面摩擦系数、空气阻力系数等全局环境变量，从而很好地重建了真实自行车赛道，并体现了道路弯道、风或天气条件等重要特征。
\item 我们在自行车与道路骑行的模型分析中采用了广为人知的牛顿力学模型，该模型简洁，并能给出科学可靠的力学分析。
\item 在严格地形式化定义优化问题之后，我们使用动态规划算法求解该优化问题。在我们的模型定义下，动态规划能够保证在相对适中的时间内求出最优解。
\end{itemize}

\subsection{缺点}

\begin{itemize}
\item 现实中的数据是连续的，而计算机仿真难以处理连续数据，因此我们对数据进行了离散化，这使连续数据损失了部分信息。
\item 影响骑手的因素还有很多，例如骑手的心理状态等，但在数学建模中我们只能考虑有限的主要因素进行建模，这可能也会损失部分影响细节。
\item 在团队计时赛比赛中，我们将多人车队简化为若干个质点，这也有可能遗漏车队内部一些能够优化表现的行为。
\end{itemize}

\section*{参考文献}

\begin{enumerate}
\item[{[1]}] F. Ashtiani, V. S. M. Sreedhara, A. Vahidi, R. Hutchison, and G. Mocko, ``Experimental modeling of cyclists fatigue and recovery dynamics enabling optimal pacing in a time trial,'' in 2019 American Control Conference (ACC). IEEE, 2019, pp. 5083--5088.
\item[{[2]}] S. A. Fayazi, N. Wan, S. Lucich, A. Vahidi and G. Mocko, ``Optimal pacing in a cycling time-trial considering cyclist's fatigue dynamics,'' 2013 American Control Conference, 2013, pp. 6442--6447, doi: 10.1109/ACC.2013.6580849.
\item[{[3]}] Rehrer, N. J., et al. ``Energy intake and expenditure during a 6-day cycling stage race.'' Scandinavian Journal of Medicine Science in Sports 20.4 (2010): 609--618.
\item[{[4]}] Sanders, Dajo, et al. ``Predicting high-power performance in professional cyclists.'' International journal of sports physiology and performance 12.3 (2017): 410--413.
\item[{[5]}] Ebert, Tammie R., et al. ``Power output during a professional men's road-cycling tour.'' International journal of sports physiology and performance 1.4 (2006): 324--335.
\item[{[6]}] https://en.wikipedia.org/wiki/Dynamic\_programming
\item[{[7]}] T. D. Noakes, A. St. Clair Gibon, and E. V. Lambert, ``From catastrophe to complexity: a novel model of integrative central neural regulation of effort and fatigue during exercise in humans: summary and conclusions,'' British J. of Sports Medicine, vol. 39, pp. 120--124, 2005.
\item[{[8]}] L. Ma, D. Chablat, F. Bennis, W. Zhang, and F. Guillaume, ``A new fatigue and recovery model and its ergonomics application in human simulation,'' Virtual and Physical Prototyping, pp. 123--137, 2010.
\item[{[9]}] https://ecal.berkeley.edu/files/ce191/CH05-DynamicProgramming.pdf
\item[{[10]}] http://ffden-2.phys.uaf.edu/webproj/211\_fall\_2016/Tad\_Nelson/Tad\_Nelson/page3.html
\item[{[11]}] https://earth.google.com/web/
\item[{[12]}] Ashtiani, Faraz; Sreedhara, Vijay Sarthy; Vahidi, Ardalan; Hutchison, Randolph; Mocko, Gregory. (2020). Optimal Pacing of a Cyclist in a Time Trial Based on Experimentally Calibrated Models of Fatigue and Recovery.
\item[{[13]}] Barry, Nathan, et al. ``Aerodynamic drag interactions between cyclists in a team pursuit.'' Sports Engineering 18.2 (2015): 93--103.
\item[{[14]}] https://www.youtube.com/watch?v=IsfnXSdpshY
\end{enumerate}

\section{科普推广材料：面向骑手与教练的实践指南}

\subsection{为什么需要科学地规划骑行}

随着生物学与物理学的发展，科学而有效的赛前准备与比赛规划在所有自行车项目中都变得越来越重要。若你与你的骑手不对比赛进行科学规划，你很可能会在长距离赛道的早期阶段耗尽体能，或者在短距离赛道上无法发挥出最佳水平。本文在此向你与你的运动员介绍我们的科学骑行方案。

\subsection{我们的方案如何提升骑手表现}

不同类型的骑手具有不同的身体表现与不同的功率曲线，功率曲线指的是骑手在某一特定时间长度内所能维持的最大功率的可视化表示。因此在自行车运动中，不同类型的骑手适合不同的战术。由此，确定你属于哪一种类型的骑手，将帮助你找到适合自己的战术，进而帮助你取得更好的成绩。

\subsection{识别骑手类型}

常见的骑手类型有：爬坡手（Climber）、陡坡冲刺手（Puncheur）、耐力型车手（Rouleur）、冲刺手（Sprinter）以及计时赛专家（Time Trial Specialist）。

我们以冲刺手与计时赛专家为例，向你展示不同类型骑手之间的差异。

冲刺手具有强大的爆发力，但输出不够稳定；而计时赛专家输出稳定，但爆发力不如冲刺手。因此在比较二者时，冲刺手适合专注于在比赛末段或途中冲刺点（若比赛设有途中冲刺点）争胜，而计时赛专家则适合要求持续输出的比赛，例如个人计时赛。

看到这里，你一定想知道自己属于哪一类骑手吧？你可以按如下步骤判断自己的类型：

\begin{enumerate}
\item 参加“最大努力”实验并获取实验数据。你可以与我们联系，也可以自行参加“最大努力”实验。该实验指的是：在充分休息之后，于平坦路段以全力冲刺，并记录你在该过程中的功率输出，进而得到你的“时间—功率”曲线。
\item 由“最大努力”实验数据获得模型参数。你可以与我们联系，也可以使用我们给出的公式，利用你的“时间—功率”曲线计算你在我模型中的参数。
\item 由模型参数得到功率曲线。使用我们的模型与算法，你可以利用第 2 步得到的参数轻松得到属于自己的功率曲线。
\item 与典型类型骑手的功率曲线作比较，确定自己的类型。将第 3 步得到的功率曲线与典型骑手的功率曲线相比较，你便能够确定自己的类型。是不是很简单？
\end{enumerate}

\subsection{提前采集道路信息}

在我们开始自行车比赛之前，骑手提前了解赛道十分重要。我们需要提前知道赛道上是否有大量陡坡？赛道上是否有大量弯道？我们还需要提前知道比赛当天的天气将会如何。所有这些对骑手决定如何骑行都非常重要。

我们的虚拟地图程序可以非常轻松地显示赛道，并计算与存储高差、曲率半径等赛道特征，以供后续分析使用。我们还建议你与你的骑手查看比赛当天的环境信息，例如风速与风向、是否下雨等，这会使我们的方案更加准确。

\subsection{方案实施}

在你依据上述各部分完成了赛道信息与骑手身体档案的详细建模之后，你可以直接使用我们提供的工具，直接获得该骑手在该赛道上的最优输出功率分配方案。若要落实该方案，我们向你提出以下建议：

\begin{itemize}
\item 使用我们提供的工具，在输入赛道信息与骑手生理数据之后，你可以得到骑手最优功率输出、速度以及赛道高度分布随赛道长度变化的可视化表示，并附有一幅动态的三维赛道路线信息图。这些图示使你与你的骑手能够非常容易地将功率分配可视化，并帮助他（她）理解最优解。
\item 若赛事组织者允许，你可以建议你的骑手携带功率计前往赛道并亲自骑行整条赛道，将仪器采集的数据与我们程序给出的最优解进行比较，这将帮助骑手更好地理解功率输出分配。
\end{itemize}

若你在确定骑手的功率档案或确定赛道信息方面存在困难，通常可以依据骑手平时的训练表现判断其类型。以下是我们针对不同类型骑手的功率分配的简单建议。

\begin{itemize}
\item \textbf{对于爬坡手：}长坡是你最有利的路线，你可以在长坡前后保存体力，把功率输出集中在长坡上，以获得更大的优势。
\item \textbf{对于陡坡冲刺手与冲刺手：}你们非常擅长在短时间内输出很大的功率，这是你们最大的优势，但这也往往使你们在长而平缓的赛道上不占优势。我们的建议是：你可以在赛道的前段与末段以及较短的爬坡段使用你的强项。
\item \textbf{对于耐力型车手与计时赛专家：}你们具备良好的综合素质组合，这通常使你们在平缓路段上占据优势。在可能的情况下，你们可以尽量避免在长坡或短坡上使出全力，否则可能在比赛后段因体能不足而落后。
\end{itemize}

\end{document}
